% ! Start justification here 1
% \textbf{Why excluding $FAR$ factor from the scoring objective.}
% Including a FAR-reduction term in the scoring objective is statistically
% redundant: in expectation, any such term reduces to the skip rate
% $S_r(\theta)$, which is already optimized directly. We demonstrate this
% using the normalized FAR-reduction term:
% \[ F(\theta) = \frac{\mathrm{FAR}_{\text{base}} - \mathrm{FAR}(\theta)}
% {\mathrm{FAR}_{\text{base}}}, \]
% where $\mathrm{FAR}_{\text{base}}$ and $\mathrm{FAR}(\theta)$ are the false
% alarm rates of the baseline and skip-enabled pipelines, respectively, from
% the scoring objective, as it is statistically redundant.

% Indeed, let $N_{\text{neg}}$ be the total number of ground-truth
% negative frames in $D_{\text{val}}$; $FP_{\text{base}}$ the number misclassified
% as positive by $\mathcal{M}$; $\mathrm{FAR}_{\text{base}} = FP_{\text{base}} /
% N_{\text{neg}}$; and $N_{\text{skip\_neg}}$ the number of negative frames
% skipped by $\mathcal{S}(\theta)$. Partition the skipped frames into $a$ frames
% that $\mathcal{M}$ \emph{would have} misclassified and $b$ frames it would have
% classified correctly, so that $a + b = N_{\text{skip\_neg}}$. Since skipped
% frames are assigned label $0$, they cannot be false positives, giving
% $FP(\theta) = FP_{\text{base}} - a$. Substituting yields:

% \[ F(\theta) = \frac{a}{FP_{\text{base}}}, \qquad \text{while} \qquad
% S_r(\theta) = \frac{a + b}{N_{\text{neg}}} \]

% Both terms share variable $a$, so the question is whether $\mathcal{S}$ tends
% to skip group-$a$ frames more often than group-$b$ frames.
% Since $\mathcal{S}$ operates solely on inter-frame pixel motion and has no
% access to $\mathcal{M}$'s weights or outputs, its skip decision is
% \emph{statistically independent} of the FP/TN label.
% Under this independence, the skip module samples blindly from the negative pool,
% so each sub-type appears in proportion to its share in the full pool.
% The expected count of group-$a$ frames among the skipped set is therefore:

% \[ \mathbb{E}[a]
% = N_{\text{skip\_neg}} \times \mathrm{FAR}_{\text{base}}
% = N_{\text{skip\_neg}} \times \frac{FP_{\text{base}}}{N_{\text{neg}}} \]

% Since $FP_{\text{base}}$ is a fixed constant, it passes through the expectation
% unchanged, and $F(\theta)$ evaluates to:

% \[ \mathbb{E}[F(\theta)]
% = \frac{\mathbb{E}[a]}{FP_{\text{base}}}
% = \frac{N_{\text{skip\_neg}} \times \dfrac{FP_{\text{base}}}{N_{\text{neg}}}}{FP_{\text{base}}}
% = \frac{N_{\text{skip\_neg}}}{N_{\text{neg}}}
% = S_r(\theta) \]

% In expectation, $F(\theta)$ is therefore identical to $S_r(\theta)$ and
% carries no independent information. This result generalizes: since
% $\mathbb{E}[\mathrm{FAR}(\theta)] = \mathrm{FAR}_{\text{base}}\,(1 -
% S_r(\theta))$ under the same independence assumption, any scoring term of
% the form $g(\mathrm{FAR}(\theta))$ is, in expectation, a monotone function
% of $S_r(\theta)$ alone. Including any FAR-based term alongside $S_r(\theta)$
% therefore adds no new information and would silently assign skip-rate
% behavior a combined weight. We therefore exclude $F(\theta)$ from the
% scoring objective and retain $\mathrm{FAR}(\theta)$ as an observed
% consequence metric reported in the results tables.


% ! Start justification here 2
% \textbf{General justification for excluding FAR-based terms from the
% scoring objective.}
% We show that, under a mild independence assumption, the expected false
% alarm rate of the skip-enabled system is fully determined by the skip
% rate $S_r(\theta)$ and the fixed baseline constant
% $\mathrm{FAR}_{\text{base}}$. This implies that any scoring term that
% depends on $\mathrm{FAR}(\theta)$ is, in expectation, redundant with
% $S_r(\theta)$ and therefore adds no independent optimization signal.

% By definition,
% $\mathrm{FAR}(\theta) = FP(\theta) / N_{\text{neg}}$,
% where $N_{\text{neg}}$ is the total number of ground-truth negative
% frames and $FP(\theta)$ is the number of false positives produced by the
% skip-enabled system. Since every skipped frame is assigned label $0$ and
% cannot be a false positive, $FP(\theta) = FP_{\text{base}} - a$, where
% $FP_{\text{base}}$ is the baseline false positive count and $a$ is the
% number of skipped frames that $\mathcal{M}$ would have misclassified.
% Taking expectations:

% \[ \mathbb{E}[\mathrm{FAR}(\theta)]
% = \frac{FP_{\text{base}} - \mathbb{E}[a]}{N_{\text{neg}}} \]

% The skip module $\mathcal{S}$ operates solely on inter-frame motion and
% has no access to $\mathcal{M}$'s weights or outputs. Its skip decision
% is therefore \emph{statistically independent} of whether a frame would
% be a false positive or a true negative under $\mathcal{M}$. Under this
% independence, $\mathcal{S}$ samples blindly from the negative pool, so
% the expected fraction of skipped frames that are false positives equals
% their prevalence in the full pool, $\mathrm{FAR}_{\text{base}}$:

% \[ \mathbb{E}[a]
% = N_{\text{skip\_neg}} \times \mathrm{FAR}_{\text{base}}
% = N_{\text{skip\_neg}} \times \frac{FP_{\text{base}}}{N_{\text{neg}}} \]


% Substituting $\mathbb{E}[a] = N_{\text{skip\_neg}} \times
% \mathrm{FAR}_{\text{base}}$ and splitting the fraction:

% \[ \mathbb{E}[\mathrm{FAR}(\theta)]
% = \frac{FP_{\text{base}}}{N_{\text{neg}}}
% - \frac{N_{\text{skip\_neg}}}{N_{\text{neg}}} \times
%   \frac{FP_{\text{base}}}{N_{\text{neg}}}
% = \mathrm{FAR}_{\text{base}} - S_r(\theta) \cdot \mathrm{FAR}_{\text{base}}
% = \mathrm{FAR}_{\text{base}} \left(1 - S_r(\theta)\right) \]

% since $N_{\text{skip\_neg}} / N_{\text{neg}} = S_r(\theta)$ by
% definition. Therefore:

% \[
% \boxed{\mathbb{E}[\mathrm{FAR}(\theta)]
% = \mathrm{FAR}_{\text{base}} \cdot \bigl(1 - S_r(\theta)\bigr)}
% \]

% This result has a direct consequence for scoring objective design.
% For any function $g$, the scoring term $g(\mathrm{FAR}(\theta))$
% satisfies, in expectation:

% \[ \mathbb{E}[g(\mathrm{FAR}(\theta))]
% = g\!\left(\mathrm{FAR}_{\text{base}} \cdot
% \bigl(1 - S_r(\theta)\bigr)\right)
% = h(S_r(\theta)) \]

% for some function $h$ that depends only on $S_r(\theta)$ and the fixed
% constant $\mathrm{FAR}_{\text{base}}$. Consequently, including any
% FAR-based scoring term $g(\mathrm{FAR}(\theta))$ alongside $S_r(\theta)$
% in the objective is equivalent to re-weighting $S_r(\theta)$ under a
% different scale --- it introduces no genuinely independent optimization
% dimension. We therefore exclude all FAR-based terms from the scoring
% objective and retain $\mathrm{FAR}(\theta)$ solely as an observed
% consequence metric reported in the results tables.


% ! % ! Start justification here 3
% \textbf{Theoretical justification for excluding FAR from the scoring objective.}
% One might consider including a normalized FAR-reduction term

% \[ F(\theta) = \frac{\mathrm{FAR}_{\text{base}} -
% \mathrm{FAR}(\theta)}{\mathrm{FAR}_{\text{base}}} \]

% which measures the relative reduction in false alarm rate with respect to the
% baseline, in the scoring objective. We argue that doing so is both structurally
% unsound and statistically redundant, and therefore exclude it.

% \textit{Structural argument.}
% The skip module $\mathcal{S}$ labels every skipped frame as negative
% ($\hat{y}_t = 0$). Since a false positive requires the system to output $1$ on a
% ground-truth negative frame, a skipped frame can never generate a new false
% positive. Consequently, $\mathrm{FAR}(\theta) \leq \mathrm{FAR}_{\text{base}}$
% holds for all $\theta$, which implies $F(\theta) \geq 0$ unconditionally.
% A scoring term that is non-negative for every feasible candidate provides no
% discriminating signal: it uniformly rewards all candidates and cannot distinguish
% configurations that skip more effectively from those that do not. FAR reduction
% is therefore a guaranteed structural side-effect of the skip mechanism, not an
% independently optimizable objective.

% \textit{Statistical argument.}
% We show that $F(\theta)$ is statistically redundant with respect to $S_r(\theta)$,
% meaning they carry the same information in expectation, and including both in the
% scoring objective is therefore unnecessary.

% Let $N_{\text{neg}}$ denote the total number of ground-truth negative
% (background-only) frames in the validation set $D_{\text{val}}$, and let
% $FP_{\text{base}}$ denote the number of those frames that the baseline classifier
% $\mathcal{M}$ incorrectly labels as fire or smoke --- its false positives.
% The baseline false alarm rate is then
% $\mathrm{FAR}_{\text{base}} = FP_{\text{base}} / N_{\text{neg}}$,
% which is the fraction of all negative frames that the classifier gets wrong.

% Now consider a candidate skip configuration $\theta$.
% Let $N_{\text{skip\_neg}}$ denote the number of negative frames that
% $\mathcal{S}(\theta)$ skips --- that is, frames assigned $s_t = 0$ without
% invoking $\mathcal{M}$.
% Among these $N_{\text{skip\_neg}}$ skipped frames, let $a$ denote the count of
% frames that $\mathcal{M}$ \emph{would have} labeled as false positives had they
% not been skipped, and let $b$ denote the count of frames it would have labeled
% correctly as true negatives, so that $a + b = N_{\text{skip\_neg}}$.

% Since every skipped frame is assigned label $0$ and therefore cannot be a false
% positive, the false positive count of the skip-enabled system satisfies
% $FP(\theta) = FP_{\text{base}} - a$, which gives
% $\mathrm{FAR}_{\text{base}} - \mathrm{FAR}(\theta)
% = (FP_{\text{base}} - FP(\theta)) / N_{\text{neg}}
% = a / N_{\text{neg}}$.
% Substituting into the definition of $F(\theta)$ yields:

% \[ F(\theta)
% = \frac{\mathrm{FAR}_{\text{base}} - \mathrm{FAR}(\theta)}{\mathrm{FAR}_{\text{base}}}
% = \frac{a / N_{\text{neg}}}{FP_{\text{base}} / N_{\text{neg}}}
% = \frac{a}{FP_{\text{base}}} \]

% Using the same counts, the skip rate $S_r(\theta)$ --- the fraction of all
% negative frames that were skipped --- expands as:

% \[ S_r(\theta) = \frac{N_{\text{skip\_neg}}}{N_{\text{neg}}} = \frac{a + b}{N_{\text{neg}}} \]

% The two terms therefore share variable $a$ but differ in what they divide by:
% $S_r$ counts all skipped negative frames $(a + b)$ over the full negative pool
% $N_{\text{neg}}$, while $F$ counts only the prevented false positives $(a)$ over
% the baseline FP count $FP_{\text{base}}$.
% The central question is whether the skip module has any systematic tendency to
% target frames in group $a$ over group $b$.

% The skip module $\mathcal{S}$ makes its decision based solely on inter-frame
% pixel motion: it compares the current frame against the previous one and skips
% if no significant motion is detected.
% It has no access to $\mathcal{M}$'s internal weights, learned feature
% representations, or prediction outputs.
% The classifier's false positives arise from background scenes that
% \emph{visually resemble} fire or smoke to $\mathcal{M}$ --- a property encoded
% entirely within the classifier's learned parameters, which is invisible to a
% motion detector.
% Therefore, the skip decision is \emph{statistically independent} of the FP/TN
% label: knowing that a frame was skipped provides no information about whether
% $\mathcal{M}$ would have misclassified it, and vice versa.

% Because of this independence, the skip module draws the $N_{\text{skip\_neg}}$
% skipped frames from the negative pool without any preference for one label over
% the other --- it samples blindly with respect to the FP/TN property.
% When items are sampled blindly from a pool, each sub-type appears in the sample
% in proportion to its share in the full pool.
% The proportion of false positives in the full negative pool is exactly
% $\mathrm{FAR}_{\text{base}} = FP_{\text{base}} / N_{\text{neg}}$.
% Applying this proportional sampling rule, the expected value
% $\mathbb{E}[\,\cdot\,]$ --- the long-run average count over many random
% outcomes --- of $a$ is:

% \[ \mathbb{E}[a] = N_{\text{skip\_neg}} \times \mathrm{FAR}_{\text{base}}
% = N_{\text{skip\_neg}} \times \frac{FP_{\text{base}}}{N_{\text{neg}}} \]

% We now compute the expected value of $F(\theta)$.
% Since $FP_{\text{base}}$ is a fixed constant determined by the baseline
% evaluation and does not vary randomly, it passes through the expectation
% unchanged:

% \[ \mathbb{E}[F(\theta)]
% = \frac{\mathbb{E}[a]}{FP_{\text{base}}}
% = \frac{N_{\text{skip\_neg}} \times \dfrac{FP_{\text{base}}}{N_{\text{neg}}}}{FP_{\text{base}}}
% = \frac{N_{\text{skip\_neg}}}{N_{\text{neg}}}
% = S_r(\theta) \]

% The $FP_{\text{base}}$ terms cancel exactly, and the result is the skip rate
% $S_r(\theta)$. Therefore, $F(\theta)$ carries no additional information beyond
% what $S_r(\theta)$ already captures: any difference between them in a specific
% experimental run is random variation, not a systematic signal attributable to
% the hyperparameter configuration $\theta$. Including both terms in the scoring
% objective would implicitly assign skip-rate behaviour a combined weight of
% $w_S + w_F$, while creating the appearance of a balanced objective with three
% independent criteria. We therefore exclude $F(\theta)$ from the scoring
% criterion and retain $\mathrm{FAR}(\theta)$ solely as an observed consequence
% metric reported in the results tables.